\documentclass[11pt]{article}
\usepackage[T1]{fontenc}
\usepackage{amsmath,amssymb,amsthm}
\usepackage[margin=1in]{geometry}
\usepackage{booktabs}
\usepackage{microtype}
\usepackage[hidelinks]{hyperref}
\hypersetup{
  pdftitle={A Proof-Carrying Exclusion Through Length 18 for the 3-out-of-5 Picture-Hanging Word},
  pdfauthor={Xiaojian Zeng}
}

\newtheorem{theorem}{Theorem}
\newtheorem{lemma}[theorem]{Lemma}
\newtheorem{proposition}[theorem]{Proposition}

\title{A Proof-Carrying Exclusion Through Length 18\\
for the 3-out-of-5 Picture-Hanging Word}
\author{Xiaojian Zeng\\
PBC School of Finance, Tsinghua University, 100190\\
\texttt{xiaojianzeng@foxmail.com}}
\date{}

\begin{document}
\maketitle

\begin{abstract}
The exact shortest-word problem for the 3-out-of-5 picture-hanging puzzle is
not settled here.  We prove the subsidiary lower bound that no admissible
element of the free group on five generators has reduced length at most 18.
The reduction has twelve finite cases: every generator has a positive even
occurrence count, a shortest word is cyclically reduced, and complete
generator, sign, and rotation symmetries leave twelve sorted count patterns
through length 18.  For each pattern, an exact interval grammar for free
reduction is translated to CNF.  Glucose3 produces a DRUP refutation, and
each archived refutation is independently accepted by \texttt{drat-trim} and
by a second RUP replay.  An explicit word of length 26 passes all twenty
essential deletion tests.  Consequently \(20\leq L(3,5)\leq 26\); the
possibilities 20, 22, 24, and 26 remain.
\end{abstract}

\section{Problem, context, and contribution}

Picture-hanging puzzles were formulated in free-group language by Demaine
et al.~\cite{DemaineEtAl2014}; later construction bounds include those of
W\"astlund~\cite{Wastlund2021}.  Verhoeff~\cite{Verhoeff2026} studied short
\(k\)-out-of-\(n\) words, proved the exact minimum 16 for 2-out-of-4, and
identified 3-out-of-5 as the next small exact-minimum case left undetermined
in that work.  The associated exhaustive-search record excludes lengths
through 16, but supplies a search log rather than a proof object for those
exclusions.

Let
\[
 F_5=\langle x_1,x_2,x_3,x_4,x_5\rangle
\]
be the free group, and let \(\ell(w)\) be the length of the unique freely
reduced word representing \(w\).  For \(S\subseteq[5]\), the deletion
homomorphism \(\pi_S\) sends \(x_i\) to \(\mathbf{1}\) for \(i\in S\) and fixes the
other generators.  An element \(w\) is \emph{3-out-of-5 admissible} when
\[
 w\ne\mathbf{1},\qquad
 \pi_S(w)=\mathbf{1}\quad(|S|=3),\qquad
 \pi_S(w)\ne\mathbf{1}\quad(|S|=2).
\]
The target of the exact problem is
\[
 L(3,5)=\min\left\{\ell(w):w\text{ is 3-out-of-5 admissible}\right\}.
\]
Equivalently, if \(\rho_R=\pi_{[5]\setminus R}\) retains precisely the
generators in \(R\), then all ten pair projections \(\rho_R(w)\),
\(|R|=2\), must be the identity and all ten triple projections, \(|R|=3\),
must be nonidentity.

Our scope is the following strict subsidiary result.

\begin{theorem}\label{thm:lower}
There is no 3-out-of-5 admissible element of reduced length at most \(18\).
Hence \(L(3,5)\geq20\).
\end{theorem}

Together with the explicit upper-bound certificate in
Proposition~\ref{prop:upper}, this narrows the exact problem to four even
values.  It does not determine which of them occurs.  Proof-carrying finite
searches elsewhere likewise separate an encoding from its solver outcome;
for example, Zeng~\cite{Zeng2026GCD} uses an exact cardinality circuit and a
checkable DRAT derivation for a finite pairwise-GCD extremal value.  Here the
domain-specific ingredient is an interval grammar that is exactly equivalent
to free reduction after each generator projection.

\section{An explicit upper-bound certificate}

Write the signed integer \(i\) for \(x_i\) and \(-i\) for \(x_i^{-1}\).
Consider
\begin{align*}
W={}&(1,2,3,4,-1,-2,-3,-4,5,\
       2,1,4,3,-4,-3,-2,-1,\\
    &\qquad 3,4,1,2,-4,-3,-2,-1,-5).
\end{align*}

\begin{proposition}\label{prop:upper}
The word \(W\) is freely reduced, has length \(26\), and is 3-out-of-5
admissible.  In particular, \(L(3,5)\leq26\).
\end{proposition}

\begin{proof}
No adjacent entries of \(W\) are inverse, so its reduced length is 26.
A direct stack reduction gives \(\rho_R(W)=\mathbf{1}\) for each of the ten pairs
\[
12,13,14,15,23,24,25,34,35,45.
\]
For the ten retained triples, the lengths of the freely reduced residuals are
as follows:
\[
\begin{array}{c|rrrrrrrrrr}
R&123&124&125&134&135&145&234&235&245&345\\ \hline
\ell(\rho_R(W))&16&16&10&18&10&10&18&10&10&10.
\end{array}
\]
Every entry is positive.  These are exactly the twenty required projection
tests, so \(W\) is admissible.
\end{proof}

The word is also obtainable from Verhoeff's one-step construction
\cite[Proposition~15]{Verhoeff2026}, with four letters cancelled at the
junction.  The displayed word and the projection data above are independent
of that provenance.

\section{Complete reduction to twelve cases}

We next isolate all cases that a shortest word of length below 20 could
occupy.

\begin{lemma}\label{lem:counts}
In every admissible word, each of the five generators occurs a positive even
number of times, half with each sign.  Consequently every admissible length is
even and at least 10.
\end{lemma}

\begin{proof}
Suppose generator \(g\) were absent.  For two other generators \(a,b\), the
projection retaining \(a,b,g\) would equal the projection retaining \(a,b\).
The former is required to be nonidentity and the latter identity, a
contradiction.  Thus every generator occurs.  Retain any pair containing
\(g\).  Its projection is the identity, so its exponent sum in \(g\), which
is the exponent sum of the original word in \(g\), is zero.  Hence the
positive and negative occurrences of \(g\) are equally numerous.
\end{proof}

\begin{lemma}\label{lem:cyclic}
A shortest admissible word is cyclically reduced.
\end{lemma}

\begin{proof}
If its reduced spelling began with a letter \(a\) and ended with \(a^{-1}\),
it would have the form \(aua^{-1}\).  Every retained-set projection of this
word is either the corresponding conjugate of the projection of \(u\), or
just that projection when \(a\) is deleted.  Thus identity versus nonidentity
is unchanged, and \(u\) would be a shorter admissible word.
\end{proof}

Let \(c_i\) be the unsigned number of occurrences of generator \(i\).
Generator permutation puts the vector \((c_1,\ldots,c_5)\) in nondecreasing
order.  Lemma~\ref{lem:counts} then gives exactly the following patterns.

\begin{table}[ht]
\centering
\caption{All sorted positive-even occurrence patterns below length 20.}
\label{tab:patterns}
\begin{tabular}{c@{\qquad}l}
\toprule
length & sorted pattern(s)\\
\midrule
10 & \((2,2,2,2,2)\)\\
12 & \((2,2,2,2,4)\)\\
14 & \((2,2,2,2,6)\), \((2,2,2,4,4)\)\\
16 & \((2,2,2,2,8)\), \((2,2,2,4,6)\), \((2,2,4,4,4)\)\\
18 & \((2,2,2,2,10)\), \((2,2,2,4,8)\), \((2,2,2,6,6)\),\\
   & \((2,2,4,4,6)\), \((2,4,4,4,4)\)\\
\bottomrule
\end{tabular}
\end{table}

The finite encoding uses only complete symmetries.  For a shortest word, first
sort its counts.  Cyclically rotate to an occurrence of a least-count
generator, call that generator 1, and independently invert generator 1 so
that the first letter is \(+1\).  Within every equal-count block, name the
remaining generators by order of first occurrence.  Independently invert each
generator whose first occurrence is negative.  Finally, cyclic reduction
ensures that the last letter is not \(-1\).  Generator permutation,
independent inversion, and cyclic rotation preserve length and every
projection condition.  Hence these constraints retain a representative of
every possible shortest word.  Reversal symmetry is not used.

\section{An exact CNF encoding of free reduction}

Fix a length \(L\) and one pattern in Table~\ref{tab:patterns}.  At each
position \(i\in\{0,\ldots,L-1\}\), one-hot variables choose one of the ten
signed generators.  Clauses forbid adjacent inverse letters.  For generator
\(g\), standard sequential-counter clauses~\cite{Sinz2005} require exactly
\(c_g/2\) positive and \(c_g/2\) negative occurrences.  The complete symmetry
conditions from the preceding section are imposed syntactically.

It remains to encode all projected free reductions.  For each retained set
\(R\) of size two or three and each half-open interval \(\mathopen{[}i,j)\), introduce
\(E_R(i,j)\).  Its intended meaning is that the projection onto \(R\) of the
subword in \(\mathopen{[}i,j)\) freely reduces to the identity.  Set \(E_R(i,i)\) true
and, in increasing interval width, impose the equivalence
\begin{equation}\label{eq:grammar}
E_R(i,j)\ \Longleftrightarrow\
\begin{cases}
E_R(i+1,j),& |w_i|\notin R,\\[2mm]
\displaystyle\mathop{\vee}\limits_{i<k<j}
 \bigl(w_k=w_i^{-1}\ \wedge E_R(i+1,k)\wedge E_R(k+1,j)\bigr),
 & |w_i|\in R.
\end{cases}
\end{equation}
All relevance tests, inverse matches, conjunctions, and disjunctions in
\eqref{eq:grammar} are represented by definitional auxiliary variables and
both directions of their CNF equivalences.  The top-level clauses require
\[
 E_R(0,L)\quad (|R|=2),
 \qquad
 \lnot E_R(0,L)\quad (|R|=3).
\]

\begin{lemma}\label{lem:grammar}
For every signed word, retained set \(R\), and interval \(\mathopen{[}i,j)\),
\(E_R(i,j)\) as defined by \eqref{eq:grammar} is true if and only if the
projected interval freely reduces to the identity.
\end{lemma}

\begin{proof}
Induct on \(j-i\).  An unretained first letter is deleted, giving the first
case.  If the first letter is retained and the projection reduces to the
identity, that letter cancels with a later inverse.  Immediately before this
cancellation the projected subword between the two letters has reduced to the
identity; after the cancellation the projected suffix must also reduce to the
identity.  This supplies one disjunct in the second case.  Conversely, any
such disjunct permits the middle and suffix reductions and then cancellation
of the matched inverse pair.  The shorter intervals are covered by the
induction hypothesis.
\end{proof}

It follows that a satisfying assignment spells a normalized, freely reduced
word with precisely the required pair and triple projections.  Conversely,
the normalized representative of any shortest admissible word with the fixed
pattern extends to a satisfying assignment by the displayed equivalences.

\section{Proof certificates and the lower bound}

The twelve formulas and archived refutations were generated under Python
3.10.12 with PySAT 1.8.dev24~\cite{IgnatievEtAl2018}, using its
\texttt{glucose3} backend.  Glucose3 returned UNSAT and emitted the textual
DRUP derivations.  A fresh run in the same recorded environment reproduced
all twelve CNFs and all twelve derivations byte for byte.  Table
\ref{tab:certificates} records the exact formula sizes and the number of
nondeletion RUP additions checked by the second replay path.  Every proof ends
with an empty-clause addition.

\begin{table}[ht]
\centering
\small
\caption{The twelve checked UNSAT cases.  ``Clauses'' counts the original
CNF; ``RUP'' counts checked clause additions, including the empty clause.}
\label{tab:certificates}
\begin{tabular}{r l r r r}
\toprule
\(L\)&pattern&variables&clauses&RUP\\
\midrule
10&(2,2,2,2,2)&6,410&31,322&35\\
12&(2,2,2,2,4)&10,228&49,810&394\\
14&(2,2,2,2,6)&15,236&74,236&759\\
14&(2,2,2,4,4)&15,270&74,278&1,733\\
16&(2,2,2,2,8)&21,612&105,398&1,465\\
16&(2,2,2,4,6)&21,658&105,428&6,117\\
16&(2,2,4,4,4)&21,696&105,506&12,613\\
18&(2,2,2,2,10)&29,516&144,096&2,838\\
18&(2,2,2,4,8)&29,574&144,142&6,264\\
18&(2,2,2,6,6)&29,582&144,194&6,700\\
18&(2,2,4,4,6)&29,624&144,208&50,966\\
18&(2,4,4,4,4)&29,666&144,294&37,708\\
\bottomrule
\end{tabular}
\end{table}

The archived CNF/proof pairs were independently replayed with
\texttt{drat-trim}~\cite{WetzlerHeuleHunt2014}.  Each of the twelve processes
returned exit code zero and exactly one \texttt{s VERIFIED} marker.  The audit
binds every result to the formula, proof, checker, and raw-output digests.

As a separate replay path, a checker uses MiniSat
2.2~\cite{EenSorensson2003} solely as a unit-propagation engine.  For every
proposed addition \(C\), it assumes the negation of every literal of \(C\),
demands a propagation conflict, and only then adds \(C\).  Thus each accepted
clause is a reverse-unit-propagation (RUP) consequence.  Deletion records are
ignored: retaining a previously checked consequence merely strengthens the
working formula, and induction shows that every retained clause follows from
the original CNF.  A RUP-checked empty clause proves unsatisfiability.  This
is the RUP fragment of the clausal proof framework discussed by Wetzler,
Heule, and Hunt~\cite{WetzlerHeuleHunt2014}.

Several checks address the finite reduction separately from solver replay.
The archived generator reproduced every stored CNF clause for clause.  An
independent partition enumeration returned respectively \(1,1,2,3,5\)
patterns at lengths \(10,12,14,16,18\).  The interval grammar was compared
with an independent stack reducer on all 335,923 words of length at most 7
over two retained generators and one ignored generator; every result agreed.
The sequential-counter component was separately tested at every length and
bound used here.  These checks are not substitutes for
Lemma~\ref{lem:grammar} or the DRUP derivations; they are implementation
audits of those two decisive components.

\begin{proof}[Proof of Theorem~\ref{thm:lower}]
Assume that an admissible element of length at most 18 exists, and choose a
shortest admissible word.  By Lemmas~\ref{lem:counts} and
\ref{lem:cyclic}, its length and sorted count pattern occur in
Table~\ref{tab:patterns}, and the word is cyclically reduced.  Apply the
complete symmetries described above.  Lemma~\ref{lem:grammar} then extends the
normalized word to a satisfying assignment of the corresponding CNF.  But
the associated checked DRUP certificate proves that CNF unsatisfiable, a
contradiction.  Therefore no admissible word has length at most 18.  Since
admissible lengths are even, \(L(3,5)\geq20\).
\end{proof}

\section{Scope, reproducibility, and next frontier}

Relative to Verhoeff's reported exclusions through length 16
\cite{Verhoeff2026}, the contribution is a proof-checkable exclusion at
length 18 together with self-contained certificates for all earlier possible
lengths.  Proposition~\ref{prop:upper} supplies the other side of
\[
20\leq L(3,5)\leq26.
\]
No assertion is made about the exact value, the existence of a length-20,
22, or 24 word, or the classification of minimizers.  The literature
comparison is limited to the cited sources and version 1 of Verhoeff's
preprint; it is not a claim of exhaustive priority.

The supplementary archive contains the CNF generator
\texttt{search\_exact\_sat.py}, all twelve DIMACS formulas and DRUP proofs,
the Glucose3 byte-reproduction audit, the independent \texttt{drat-trim}
audit, the separate RUP replay checker \texttt{check\_drup.py}, the
deterministic encoding audit \texttt{audit\_exact\_encoding.py}, and the
independent upper-word checker \texttt{verify\_upper26.py}.  The next
decisive computation is the same proof-producing split at length 20.  Its
seven sorted patterns are
\[
\begin{gathered}
(2,2,2,2,12),\ (2,2,2,4,10),\ (2,2,2,6,8),\
(2,2,4,4,8),\\
(2,2,4,6,6),\ (2,4,4,4,6),\ (4,4,4,4,4).
\end{gathered}
\]
A checked satisfying word in any case would prove \(L(3,5)=20\); seven
checked UNSAT certificates would raise the lower bound to 22.

\paragraph{Computational and AI assistance.}
The finite search, proof generation, byte reproduction, and exact replay used
the software and versions stated above; no floating-point calculation
supports a mathematical claim.  An OpenAI Codex language model assisted with
research planning, implementation, manuscript drafting, and internal
consistency checks.  The author reviewed the statement, proof, code, and
reported scope and takes responsibility for the article.

\bibliographystyle{plain}
\bibliography{references}

\end{document}
